\documentclass[11pt]{article}

\usepackage[margin=1in]{geometry}
\usepackage{amsmath,amssymb,amsthm}
\usepackage{booktabs}
\usepackage[round]{natbib}
\usepackage[hidelinks]{hyperref}
\usepackage{microtype}

\newtheorem{proposition}{Proposition}
\newtheorem{corollary}{Corollary}
\newtheorem{lemma}{Lemma}
\theoremstyle{definition}
\newtheorem{definition}{Definition}
\newtheorem{remark}{Remark}

\newcommand{\R}{\mathbb{R}}
\newcommand{\E}{\mathbb{E}}
\renewcommand{\Pr}{\mathbb{P}}
\newcommand{\KL}{\mathrm{KL}}
\newcommand{\given}{\,\vert\,}
\DeclareMathOperator*{\argmax}{arg\,max}

\title{\textbf{Marginally Useful}\\[2pt]
\large Conformal Prediction Certifies the Coverage of a Set, \\ Not the Quality of a Distribution}

\author{Peter Cotton\thanks{Microprediction. \texttt{peter.cotton@microprediction.com}.
This paper accompanies the interactive demonstrations at \emph{Marginally Useful};
every quantitative claim below has a slider-driven counterpart there.}}

\date{\today}

\begin{document}
\maketitle

\begin{abstract}
Conformal prediction is correct mathematics with a genuine guarantee, and it is
routinely filed in the wrong drawer. The guarantee it delivers --- finite-sample,
distribution-free \emph{marginal} coverage of a set --- is not the guarantee that
forecasting, density estimation, and decision-making usually require. This paper is
mostly expository: we assemble, in one place and with citations, the reasons the
marginal guarantee is weaker than it looks (it is provably not conditional; it is
trivially satisfiable; it evaporates without exchangeability). We then offer three
small contributions, two of which are folklore made precise and one of which we
believe is a useful new lens. First, a constructive \emph{orthogonality} proposition:
coverage and predictive log-score live on independent axes, in the strong sense that
one can hold the conformal set (hence marginal coverage) exactly fixed while moving
the log-score by an arbitrary amount. Second, a \emph{sharpness-gap} decomposition:
the expected log-score regret of a conformal predictive system relative to the oracle
equals the expected Kullback--Leibler divergence of the conditional residual law from
its marginal --- a quantity that is structurally invariant to \emph{any} marginal
recalibration, conformal included, and can be reduced only by conditioning on the
input. Within the class of forecasters it actually inhabits, the conformal predictive
system is log-score optimal; the cost is the class, not the calibration. Third, the
\emph{coverage--sharpness plane}, a diagnostic in which conformalization is exactly a
horizontal projection. We close with the honest converse: a precise litmus test for
when coverage \emph{is} the objective, and conformal prediction is exactly right.
\end{abstract}

\section{Introduction}

A practitioner maintains a probabilistic forecaster and is told, repeatedly, that
conformal prediction is the principled, distribution-free way to ``add uncertainty.''
They bolt it on. The number they actually report --- a proper score such as the
predictive log-likelihood --- does not improve. The natural diagnosis is a tuning
failure. The correct diagnosis, we argue, is a \emph{category error}: conformal
prediction is a method for \emph{certifying the coverage of a set}, and it is being
asked to \emph{estimate a distribution}. These are different kinds of object, graded
by different instruments, and no amount of tuning converts one into the other.

None of this impugns the mathematics. Given exchangeable data, a fitted predictor,
and a nonconformity score, split conformal prediction returns a set $C_\alpha(x)$ with
\begin{equation}\label{eq:marginal}
  \Pr\!\big(Y_{n+1}\in C_\alpha(X_{n+1})\big)\;\ge\;1-\alpha,
\end{equation}
finite-sample and with no assumption on the data distribution
\citep{vovk2005algorithmic,lei2018distribution,angelopoulos2023gentle}. The trouble
is entirely in the reading of \eqref{eq:marginal}. The probability averages over
$X_{n+1}$ as well as $Y_{n+1}$: it is a statement about the \emph{average} future
input, not about the one in front of you. It is \emph{marginal} coverage, and as the
foundational analyses warned from the start, ``a good estimator must satisfy something
more than marginal coverage'' \citep{lei2014distribution}.

\paragraph{Contributions and structure.}
Sections~\ref{sec:setup}--\ref{sec:marginal} are expository: the guarantee, and the
three senses in which it is weaker than advertised. Our contributions are:
\begin{itemize}
  \item \textbf{Orthogonality} (Section~\ref{sec:orthogonality},
  Proposition~\ref{prop:orth}). A constructive statement that the conformal set is a
  function of the nonconformity scores alone, so marginal coverage is statistically
  independent of every feature of the predictive distribution the score does not
  encode --- sharpness in particular. One can pin coverage at $1-\alpha$ and send the
  log-score to $-\infty$. \emph{(Folklore made precise.)}
  \item \textbf{The sharpness gap} (Section~\ref{sec:gap},
  Propositions~\ref{prop:relevel}--\ref{prop:gap}). The conformal predictive system
  re-levels the base model's residual shape; its log-score regret to the oracle is
  exactly $\E_X\,\KL\!\big(r(\cdot\given X)\,\|\,\bar r\big)$, the expected divergence
  of the conditional residual law from its marginal. This term is invariant to any
  marginal recalibration and is the precise content of ``blind to sharpness.''
  Within the single-shape class, the conformal system is log-score optimal.
  \emph{(One known identity restated; one decomposition we have not seen stated this
  way.)}
  \item \textbf{The coverage--sharpness plane} (Section~\ref{sec:plane}). A two-axis
  diagnostic in which conformalization is a horizontal projection onto the
  zero-coverage-error line, leaving the vertical (score) coordinate untouched.
  \emph{(A presentation device; the lens is, we think, clarifying.)}
\end{itemize}
Section~\ref{sec:exch} treats exchangeability and time series; Section~\ref{sec:steel}
gives the honest converse --- when coverage really is the objective.

\section{Setup: split conformal and the marginal guarantee}\label{sec:setup}

Let $(X_i,Y_i)_{i=1}^{n+1}$ be exchangeable. Fix a point predictor $\hat\mu$ trained on
a disjoint set, and a \emph{nonconformity score} $s(x,y)$; the canonical choice in
regression is the absolute residual $s(x,y)=|y-\hat\mu(x)|$. Compute calibration
scores $s_i=s(X_i,Y_i)$ for $i=1,\dots,n$ and let
\begin{equation}\label{eq:q}
  \hat q \;=\; s_{(k)}, \qquad k=\big\lceil (n+1)(1-\alpha)\big\rceil,
\end{equation}
the $k$-th order statistic (with $\hat q=+\infty$ when $k>n$). The split-conformal set is
\begin{equation}\label{eq:set}
  C_\alpha(x)=\{y:\ s(x,y)\le \hat q\}
  \;=\;[\hat\mu(x)-\hat q,\ \hat\mu(x)+\hat q]
\end{equation}
for the absolute-residual score. Exchangeability of the augmented scores makes the rank
of $s_{n+1}$ uniform, which yields \eqref{eq:marginal}
\citep{vovk2005algorithmic,lei2018distribution}. The construction is elegant and the
guarantee is real; the demonstrations confirm that the empirical coverage tracks
$1-\alpha$ as the sample size, noise level, and $\alpha$ are varied. We do not dispute
any of this. We dispute what it is taken to mean.

\section{Three ways the marginal guarantee is weaker than it looks}\label{sec:marginal}

\paragraph{(i) Marginal is not conditional, and conditional is impossible for free.}
One would prefer \emph{conditional} coverage,
$\Pr\big(Y\in C(x)\given X=x\big)\ge 1-\alpha$ for (almost) every $x$. Distribution-free,
this cannot be bought. \citet{lei2014distribution} show (their Lemma~1) that any band
with non-trivial finite-sample conditional validity has \emph{infinite} expected length
at almost every point of a continuous distribution.
\citet{barber2021limits} restate and sharpen this: a conditionally valid $C_n$ has
$\E\,\mathrm{leb}(C_n(x))=\infty$ at almost all non-atomic $x$ (their Prop.~2.2), and
even the relaxed demand of coverage on every subgroup of probability $\ge\delta$ ``is
impossible to attain beyond the trivial solution'' of inflating the marginal level to
$1-\alpha\delta$ (their Thm.~3.1). We collect both as the operative no-go result.

\begin{lemma}[No-go for distribution-free conditional coverage]\label{lem:nogo}
Let $P$ have non-atomic $X$. \emph{(a)} \citep{lei2014distribution,barber2021limits}
Any $C_n$ with non-trivial finite-sample conditional validity has
$\E\,\mathrm{leb}\big(C_n(x)\big)=\infty$ at almost every $x$. \emph{(b)}
\citep{barber2021limits} Requiring $\Pr\big(Y\in C_n(X)\given X\in G\big)\ge 1-\alpha$
for every measurable $G$ with $P(X\in G)\ge\delta$ is attainable only by a $C_n$ whose
marginal level is inflated to $1-\alpha\delta$ --- a uniformly wider, non-adaptive band.
\end{lemma}

\noindent These are not loose bounds to be tightened; they are the reason the gap cannot
be closed for free, and both surface as concrete finite-sample facts in the
demonstrations. Part~(a) drives Demo~7: localizing split conformal to recover per-$x$
coverage shrinks each cell's calibration count $n_b$ until
$\lceil (n_b+1)(1-\alpha)\rceil > n_b$, at which point \eqref{eq:q} gives $\hat q=+\infty$
and the only valid cell interval is $(-\infty,\infty)$ --- the infinite expected length
of the lemma, arriving as arithmetic. Part~(b) is Demo~8: the distribution-free
protection of all $\delta$-subgroups is exactly the flat band at level $1-\alpha\delta$,
whose width diverges as $\delta\to 0$ while never adapting to $x$, at a multiplicative
width cost over an assumption-driven adaptive oracle.

So the guarantee you can have distribution-free is
the marginal one, and a marginally valid band ``tends to overestimate the set when $x$
is in the high density area and to underestimate for low density $x$''
\citep{lei2014distribution} --- it over-covers the easy inputs and under-covers the hard
ones, silent by construction about the cases the model was built to get right. The
heteroscedastic demonstration (Demo~2) exhibits exactly this: $90\%$ overall,
near-$100\%$ where the noise is small, well below target where it is large.

\paragraph{(ii) Validity is trivially satisfiable.}
A predictor that returns the whole outcome space with probability $1-\alpha$ and the
empty set otherwise satisfies \eqref{eq:marginal} exactly and conveys nothing. More
practically, the \emph{same} coverage attaches to an excellent $\hat\mu$ and to a
deliberately terrible one --- the latter merely yields a wider $\hat q$ in
\eqref{eq:q}. Marginal validity is therefore not evidence of a good method; it
certifies the fence, not the cattle. The next section makes this precise.

\paragraph{(iii) Without exchangeability, even the marginal guarantee fails.}
Equation~\eqref{eq:marginal} rests on exchangeability of the augmented sample. Under
covariate shift, label shift, or temporal dependence it does not hold, and the patches
that restore something do so by weakening what is guaranteed (Section~\ref{sec:exch}).

\section{Coverage and score are orthogonal}\label{sec:orthogonality}

We first make ``coverage is not a measure of forecast quality'' a precise, constructive
fact. Model a forecaster as a pair $(\,f(\cdot\given x),\,s\,)$ of a predictive density
and a nonconformity score.

\begin{proposition}[Orthogonality of coverage and log-score]\label{prop:orth}
The split-conformal set \eqref{eq:set} is a measurable function of the nonconformity
scores $\{s_i\}$ and of $s(x,\cdot)$ alone. Consequently, if two forecasters share the
same score function $s$ and the same calibration scores, their conformal sets --- and
hence their marginal coverage at every level $\alpha$ --- are \emph{identical}, while
their expected log-scores $\E[\log f(Y\given X)]$ may differ by an arbitrary amount.
\end{proposition}

\begin{proof}
The first claim is immediate from \eqref{eq:q}--\eqref{eq:set}: $\hat q$ is an order
statistic of $\{s_i\}$ and $C_\alpha(x)=\{y:s(x,y)\le\hat q\}$, so two forecasters
agreeing on $s$ and on $\{s_i\}$ produce the same set for every $x,\alpha$, and
\eqref{eq:marginal} concerns only that set.

For the gap, take $X$ degenerate and $Y\sim N(0,1)$; let $\hat\mu\equiv 0$ with the
absolute-residual score $s(x,y)=|y|$, which does not depend on the predictive density at
all. For any claimed predictive density $f_\sigma=N(0,\sigma^2)$ the scores are
$|Y_i|$, independent of $\sigma$, so the conformal set is the fixed interval
$[-\hat q,\hat q]$ with $\hat q$ the empirical $(1-\alpha)$-quantile of $|Y|$. Yet
\[
  \E[\log f_\sigma(Y)] \;=\; -\tfrac12\log(2\pi\sigma^2)-\frac{\E[Y^2]}{2\sigma^2}
  \;=\; -\tfrac12\log(2\pi\sigma^2)-\frac{1}{2\sigma^2}\;\xrightarrow[\sigma\to 0^+]{}-\infty,
\]
and likewise $\to-\infty$ as $\sigma\to\infty$. Hence for any $M>0$ there are
$\sigma_1,\sigma_2$ with identical conformal sets (identical marginal coverage at every
$\alpha$) yet log-scores differing by more than $M$.
\end{proof}

\begin{corollary}\label{cor:orth}
When $s$ is the residual $|y-\hat\mu(x)|$, marginal validity certifies a property of the
residual order statistics and is statistically independent of the scale, shape, and
sharpness of $f$. In particular, no inference about forecast quality (as graded by a
strictly proper score) can be drawn from coverage alone.
\end{corollary}

This is the rigorous form of a point made informally by several authors --- that
conformal validity ``almost invites you to use garbage prediction functions''
\citep{recht2024cover} --- and it is what the orthogonality demonstration shows
visually: the claimed density's spread slides freely while the conformal interval and
its $90\%$ coverage do not move, because the spread never entered \eqref{eq:q}.

\section{The category error, made precise: the sharpness gap}\label{sec:gap}

Forecast quality decomposes into \emph{calibration} and \emph{sharpness}, and the
discipline's stated goal is to maximize sharpness subject to calibration, assessed by a
proper scoring rule \citep{gneiting2007probabilistic,gneiting2007strictly}. We show
where conformal prediction sits in this decomposition, and we do so fairly: within the
class of objects it can produce, it is optimal; the loss is the class.

\paragraph{Conformal predictive systems.} Sweeping $\alpha$ in \eqref{eq:set} and
stacking the nested intervals yields a full predictive distribution --- the
\emph{conformal predictive system} (CPS) of \citet{vovk2019nonparametric}. This is the
fairest possible target for the criticism, because it is the one mode in which conformal
prediction outputs a distribution rather than a set.

\begin{proposition}[Re-leveling identity]\label{prop:relevel}
Let $\hat\mu$ be a fixed location predictor with residual $R=Y-\hat\mu(X)$ and
residual-based score, and let $\Pi_x$ be the CPS predictive distribution. Under
exchangeability, $\Pi_x(y)=\widehat G_n\big(y-\hat\mu(x)\big)$, where $\widehat G_n$ is
the empirical CDF of the calibration residuals; and $\widehat G_n\to G$ uniformly
(Glivenko--Cantelli), with $G$ the CDF of $R$. The CPS predictive density is therefore,
in the large-calibration limit, $\pi_x(y)=g\big(y-\hat\mu(x)\big)$ with $g$ the
\emph{marginal} residual density.
\end{proposition}

\begin{proof}[Proof sketch]
The CPS at level $\alpha$ returns the interval whose endpoints are the residual order
statistics determined by \eqref{eq:q}; as $\alpha$ varies these endpoints trace out
$\hat\mu(x)$ plus the residual quantile function, i.e.\ the inverse of $\widehat G_n$
shifted by $\hat\mu(x)$. Inverting gives $\Pi_x(y)=\widehat G_n(y-\hat\mu(x))$. The
limit is Glivenko--Cantelli. See \citet{vovk2019nonparametric} for the exact finite-sample
construction.
\end{proof}

The estimation, in other words, was done by the base model and the residual sample; the
CPS supplies only the leveling. We now quantify what that costs. Write the true
conditional residual density as $r(\cdot\given x)$, so the oracle density is
$q^\star(y\given x)=r\big(y-\hat\mu(x)\given x\big)$, and let
$\bar r(\cdot)=\E_X\,r(\cdot\given X)=g(\cdot)$ be its marginal.

\begin{proposition}[Sharpness gap]\label{prop:gap}
Fix the location predictor $\hat\mu$. Then the expected log-score regret of the CPS
predictive density $\pi_X$ relative to the oracle is
\begin{equation}\label{eq:gap}
  \E\big[\log q^\star(Y\given X)\big]-\E\big[\log \pi_X(Y)\big]
  \;=\; \E_X\,\KL\!\big(r(\cdot\given X)\ \big\|\ \bar r\big)\ \ge\ 0,
\end{equation}
with equality if and only if $R\perp X$ (the residual law does not depend on the input).
Moreover, among all forecasters of the single-shape location form
$y\mapsto h\big(y-\hat\mu(x)\big)$ for a shared density $h$, the expected log-score
$\E[\log h(R)]$ is maximized at $h=g$; the CPS attains this maximum. Consequently, any
procedure that returns a single shared residual shape --- conformal predictive systems
and \emph{marginal} recalibration alike --- leaves the gap \eqref{eq:gap} unchanged.
\end{proposition}

\begin{proof}
With $R=Y-\hat\mu(X)$ and $\pi_X(Y)=g(R)$ we have $\log q^\star(Y\given X)=\log r(R\given X)$, so
\[
  \E[\log q^\star(Y\given X)]-\E[\log \pi_X(Y)]
  =\E\Big[\log\tfrac{r(R\given X)}{g(R)}\Big]
  =\E_X\,\E_{R\given X}\Big[\log\tfrac{r(R\given X)}{g(R)}\Big]
  =\E_X\,\KL\!\big(r(\cdot\given X)\,\|\,g\big).
\]
Each inner term is a KL divergence, hence $\ge 0$, with equality iff
$r(\cdot\given X)=g$ a.s., i.e.\ $R\perp X$. For the second claim, for any shared shape
$h$, $\E[\log g(R)]-\E[\log h(R)]=\KL(g\,\|\,h)\ge 0$ by Gibbs' inequality, so $h=g$ is
optimal; by Proposition~\ref{prop:relevel} the CPS realizes $h=g$ in the limit. A
forecaster constrained to a single shape cannot represent $X$-dependence in $r$, so it
cannot alter the right-hand side of \eqref{eq:gap}.
\end{proof}

\begin{remark}[What this says, and does not say]
The gap \eqref{eq:gap} is the expected divergence of the conditional residual law from
its marginal --- a pure \emph{conditional-shape} (heteroscedasticity-and-beyond) term.
It is the exact, quantitative content of ``constitutionally blind to sharpness'': the
blindness is invariant to anything marginal. Three honest caveats. (1)
\emph{Conformal is not dominated within its class.} The proposition shows the CPS is
log-score optimal among single-shape location forecasters; the cost is the restriction,
not the calibration step. (2) \emph{Recalibration weakly dominates only by escaping the
class.} Distributional recalibration fit to a proper score \citep{kuleshov2018accurate}
can move along the shape axis when it conditions on $x$ (heteroscedastic modeling,
conditional PIT maps), and only then does it strictly reduce \eqref{eq:gap}; marginal
PIT recalibration alone does not. Conformalized quantile regression
\citep{romano2019conformalized} is the conformal world's way of buying $x$-dependent
shape, and it is no longer in the single-shape class. (3) \emph{The distribution-free
guarantee is real.} Recalibration offers calibration only in expectation or
asymptotically and carries no coverage certificate; conformal's marginal guarantee is
finite-sample and assumption-light. When the certificate itself is the deliverable, that
is decisive (Section~\ref{sec:steel}).
\end{remark}

The recalibration demonstration realizes the special case in which $\hat\mu$ is correct
and $r(\cdot\given x)$ is a fixed Gaussian shape with the wrong scale: variance
recalibration fit to the log-score recovers the right shape and improves the score,
while conformalization re-levels to exact coverage and returns a set.

\section{The coverage--sharpness plane}\label{sec:plane}

The preceding results suggest a single diagnostic. Place a forecaster at coordinates
\begin{equation}
  \big(\,C,\,S\,\big),\qquad
  C=\big|\,\Pr(Y\in C_\alpha(X))-(1-\alpha)\,\big|,\qquad
  S=\E[\log f(Y\given X)],
\end{equation}
the marginal-coverage error on the horizontal axis and the mean predictive log-score on
the vertical. The decomposition of forecast quality into calibration and sharpness
\citep{gneiting2007probabilistic} is, in this plane, the statement that the two
coordinates are distinct. Propositions~\ref{prop:orth}--\ref{prop:gap} then read
geometrically:
\begin{itemize}
  \item \textbf{Conformalization is a horizontal projection.} It drives $C\to 0$ and,
  by Proposition~\ref{prop:orth}, leaves $S$ at the base model's value. It moves a point
  left, never up.
  \item \textbf{Recalibration and better modeling move vertically.} Fitting shape to a
  proper score raises $S$; conditioning on $x$ is the only way to raise it past the
  ceiling set by the sharpness gap \eqref{eq:gap}.
\end{itemize}
A method reported by its coverage alone is reported by one coordinate of a
two-coordinate quality. The plane is a presentation device --- adjacent to reliability
and sharpness diagrams \citep{gneiting2007probabilistic} --- but the ``conformal = move
left, never up'' reading makes the category distinction hard to miss, and it is the
through-line of all six demonstrations.

\section{Exchangeability, and why time series are the worst case}\label{sec:exch}

Drop exchangeability and \eqref{eq:marginal} is no longer guaranteed. The literature's
repairs are ingenious, and uniform in one respect: they survive by redefining what is
guaranteed.
\begin{itemize}
  \item \textbf{Weighted conformal prediction} \citep{tibshirani2019conformal} restores
  marginal coverage under covariate shift when the likelihood ratio is known; with
  estimated weights coverage becomes approximate and the effective sample size drops.
  \item \textbf{Adaptive Conformal Inference} \citep{gibbs2021adaptive} updates the level
  online and ``puts no constraints on the data generating distribution,'' but guarantees
  a \emph{long-run time-average} miscoverage frequency,
  $\frac1T\sum_t \mathrm{err}_t\to\alpha$ --- not a finite-sample statement about
  $\Pr(Y\in C(X))$, and certainly not conditional coverage.
  \item \textbf{EnbPI} \citep{xu2021conformal} drops exchangeability for approximate
  marginal coverage under strong-mixing/stationarity surrogates; later online methods
  \citep{zaffran2022adaptive,gibbs2024conformal} retreat to regret bounds over local
  intervals. The general beyond-exchangeability analysis of \citet{barber2023conformal}
  bounds the coverage gap by a total-variation distance that is exactly what
  nonstationarity inflates.
\end{itemize}
The pattern is a sequence of weaker objects: from finite-sample per-instance coverage to
a long-run average frequency to a regret bound. Each is respectable; none is the sharp
predictive distribution \emph{now}, for \emph{this} step, that time-series forecasting
asks for and that a proper score grades. The time-series demonstration shows fixed split
conformal collapsing well below target once drift sets in, while the adaptive patch
restores the long-run average yet leaves any short window swinging around it --- average,
not now.

\section{When coverage \emph{is} the objective}\label{sec:steel}

The argument is skeptical, not dismissive, and it comes with a precise converse. The
litmus is one question:
\begin{quote}
\emph{Is your loss a function of \textbf{whether} the truth lands in a region, or of
\textbf{where} it lands?}
\end{quote}
If the former, coverage is the native object and conformal prediction is sound and often
ideal:
\begin{itemize}
  \item \textbf{Selective prediction and risk control.} ``Return a small set guaranteed
  to contain the label $95\%$ of the time; a human checks it.'' The deliverable is a set;
  risk-controlling prediction sets formalize this \citep{angelopoulos2021uncertainty,
  bates2021distribution}.
  \item \textbf{Retrieval and shortlisting,} where the product is a candidate set and
  coverage is recall.
  \item \textbf{Anomaly and novelty detection via conformal $p$-values} --- conformal
  prediction's cleanest home, because here it is not a forecaster but a
  distribution-free hypothesis test and coverage is Type-I error control
  \citep{vovk2005algorithmic}.
  \item \textbf{Compliance and long-run frequency control,} where a certificate of the
  form ``contained $95\%$ of the time on average'' \emph{is} the product, and the
  distribution-free, finite-sample nature of \eqref{eq:marginal} is exactly the value
  added.
\end{itemize}
Two caveats keep even these honest. If you trust a calibrated posterior you can form a
sharper set yourself --- the smallest region of mass $\ge 1-\alpha$ --- and recover the
expectations a bare set discards; conformal's marginal value is then purely the
distribution-free guarantee, useful when you distrust the density and exchangeability
holds. And even when coverage is the goal you usually want it conditionally, which
returns us to the impossibility of Section~\ref{sec:marginal}.

\section{Conclusion}

Conformal prediction certifies the coverage of a set, finite-sample and
distribution-free, under exchangeability. That is the one true thing it does, and it is
the wrong thing for forecasting: the guarantee you can have (marginal) is rarely the one
you want (conditional); the conditional one is provably unavailable distribution-free;
the output is a set, not a measure; it is blind to sharpness in the exact amount
\eqref{eq:gap}; and the exchangeability it needs fails in every time series. Within the
class of objects it produces it is optimal, and where a guarantee is the product it is
the right tool. But coverage answers a question --- \emph{is the truth in this region?}
--- that, in forecasting, you were usually not asking. Confusing the answer with an
estimate of the distribution is a category error; seen as a type error rather than a
tuning problem, the disappointment stops being mysterious.

\paragraph{Reproducibility.} Every numerical claim has an interactive counterpart in the
accompanying \emph{Marginally Useful} demonstrations; the orthogonality of
Section~\ref{sec:orthogonality}, the gap of Section~\ref{sec:gap}, and the projection of
Section~\ref{sec:plane} are each a slider away.

\small
\bibliographystyle{plainnat}
\bibliography{references}

\end{document}
